% Chapter 19
\chapter{Accessibility Relaxation}

\chapterepigraph{%
  The universe does not expand\\
  into anything.\\
  It relaxes\\
  away from tension.
}{Flyxion, \textit{Axioms for a Falling Universe}}

\sidenote{App.\,J §J.21\\on accessibility\\relaxation flows}

Standard cosmology explains the redshift of distant galaxies through
the expansion of space: the fabric of spacetime stretches, carrying
galaxies apart and stretching the wavelengths of photons traveling
between them. This chapter develops the RSVP alternative: \emph{accessibility
relaxation}, in which redshift arises not from spatial expansion but
from the progressive decrease of accessibility entropy along cosmological
light paths.

\section{Replacing Expansion}

\begin{tcolorbox}[enhanced, breakable,
  colback=RSVPwarm, colframe=RSVPaccent, arc=4pt, boxrule=0.8pt,
  title={\bfseries\sffamily Worked Example: Two Falsifiable Predictions of RSVP Cosmology},
  coltitle=white, attach boxed title to top left={yshift=-2mm,xshift=6mm},
  boxed title style={colback=RSVPaccent,sharp corners},
  top=4mm,left=4mm,right=4mm,bottom=4mm]

The RSVP framework is not merely a reinterpretation of existing observations —
it makes specific predictions that differ numerically from $\Lambda$CDM.

\textbf{Prediction 1: The third-order redshift-luminosity coefficient.}
The coefficient $c_3$ in the Taylor expansion of the redshift-luminosity
relation measures how cosmic ``acceleration'' varies with distance.
\begin{itemize}
  \item $\Lambda$CDM predicts: $c_3 \approx -0.1804$\quad (negative)
  \item RSVP predicts: $c_3 \approx +2.3591 + \delta_{\Phi,v}$\quad (positive)
\end{itemize}
where $\delta_{\Phi,v}$ is a small correction from the $\Phi$-$\mathbf{v}$ coupling.
The sign difference is decisive: point telescopes at deep-field supernovae,
measure $c_3$, and one of the two frameworks is immediately falsified.

\textbf{Prediction 2: Non-zero redshift-birefringence cross-correlation.}
In an empty-vacuum cosmology, the stretching of light (redshift) and the
twisting of light's polarization (birefringence) are independent phenomena.
RSVP predicts they are correlated, because both arise from the same
underlying entropy-vorticity coupling in the lamphrodyne flow field $\mathbf{v}$:
\[
  \langle \delta z(\hat{n})\,\delta\alpha(\hat{n}')\rangle \neq 0
\]
where $\delta z$ is the redshift perturbation and $\delta\alpha$ is the
birefringence angle. This cross-correlation is zero by symmetry in $\Lambda$CDM.
Missions such as \textit{Euclid} or the \textit{Nancy Grace Roman Space Telescope}
will measure this directly.

If either prediction holds, the relaxing-spring picture is confirmed.
If both fail, the RSVP framework requires fundamental revision.
\end{tcolorbox}


The Hubble relation $v = H_0 d$ (recession velocity proportional to
distance) is observational. Its \emph{interpretation} as spatial expansion
is theoretical. RSVP proposes a different interpretation: the Hubble
relation reflects a gradient in the accessibility entropy field $S$,
not an expansion of space.

\begin{definition}{RSVP Hubble Relation}{rsvp-hubble}
  In the RSVP framework, the observed redshift of a photon traveling
  from source to observer along path $\gamma$ is:
  \[
    z = \int_\gamma \alpha \cdot |\nabla S| \,dl
  \]
  where $\alpha$ is a coupling constant and $|\nabla S|$ is the magnitude
  of the accessibility gradient along the photon path.
\end{definition}

This replaces the geometric interpretation (space is stretched) with a
field-theoretic one (photons lose energy as they traverse accessibility
gradients). The two interpretations are empirically distinguishable:
the RSVP interpretation predicts a specific relationship between redshift,
cosmic structure (which determines $|\nabla S|$), and direction (accessibility
gradients are not isotropic on small scales).

\section{Accessibility Gradients}

The accessibility gradient $\nabla S$ is the key dynamical object
in RSVP cosmology. It determines:

\begin{itemize}
  \item The direction of cosmological flow (matter flows toward
        accessibility sinks).
  \item The magnitude of cosmological redshift (photons lose energy
        proportional to the gradient they traverse).
  \item The formation of large-scale structure (filaments and voids
        form along accessibility ridges and wells).
\end{itemize}

\begin{definition}{Accessibility Gradient Cosmology}{acc-grad-cosmo}
  An \emph{accessibility gradient cosmology} is a RSVP model in which
  the large-scale behavior of the universe is determined by the gradient
  field $\nabla S$ rather than by a metric expansion parameter $a(t)$.
\end{definition}

\section{Entropic Smoothing}

Over cosmological time, the accessibility entropy field $S$ tends to
smooth: sharp gradients diffuse, accessibility wells fill, and the
landscape approaches a flatter configuration. This is \emph{entropic
smoothing} — the RSVP version of the heat death of the universe.

\begin{definition}{Entropic Smoothing}{entropic-smoothing}
  \emph{Entropic smoothing} is the long-time behavior of the accessibility
  field under the diffusion term $D_S \nabla^2 S$ in the $S$-field equation.
  As $t \to \infty$, $S \to S_\infty$ (a constant), and $\nabla S \to 0$.
\end{definition}

The approach to $S_\infty$ represents the exhaustion of accessible futures:
all structures have formed, all accessible configurations have been reached,
all gradients have been eliminated. This is the RSVP interpretation of
the heat death: not maximum disorder, but minimum accessibility gradient —
a state with nowhere left to flow.

\section{Lamphron–Lamphrodyne Dynamics Revisited}

The lamphron–lamphrodyne distinction of Chapter~18 is central to
understanding the formation and dissolution of cosmic structures.

Lamphron regions (accessibility maxima) are sources: they drive flow
outward, pushing scalar density toward lamphrodyne sinks. As matter
accumulates at lamphrodyne sinks, it forms structures. As structures
form, the local $S$ decreases (accessibility is consumed), eventually
turning the lamphrodyne region into a lamphron as the structure radiates
energy into its environment.

This cycle — lamphron drives flow to lamphrodyne, lamphrodyne forms
structure, structure radiates and becomes lamphron — is the RSVP
account of the stellar life cycle, galactic evolution, and potentially
the life cycle of the universe as a whole.

\begin{exercise}
  Derive the entropic smoothing timescale: the characteristic time
  $\tau_S = L^2 / D_S$ for accessibility gradients of length scale $L$
  to diffuse under the equation $\partial_t S = D_S \nabla^2 S$.
  What physical observable corresponds to $D_S$ in the RSVP framework?
\end{exercise}

\begin{exercise}
  In standard cosmology, the Hubble constant $H_0$ has units of inverse
  time ($\text{km/s/Mpc}$). In the RSVP Hubble relation, the coupling
  constant $\alpha$ and the accessibility gradient $|\nabla S|$ must
  combine to give units of inverse length. What are the units of
  $|\nabla S|$ and $\alpha$ separately?
\end{exercise}
